% !TEX program = lualatex
% Transparencias de Fundamentos de las Matemáticas I
% Clase 06: Conjuntos I

\documentclass[aspectratio=169,11pt]{beamer}

\usepackage{fontspec}
\usepackage[spanish,es-nodecimaldot]{babel}
\usepackage{amsmath,amssymb,mathtools}
\usepackage{booktabs}
\usepackage{csquotes}
\usepackage{xcolor}

\usetheme{Madrid}
\usecolortheme{default}
\setbeamertemplate{navigation symbols}{}
\setbeamertemplate{footline}[frame number]

\definecolor{FdMBlue}{RGB}{20,55,100}
\definecolor{FdMRed}{RGB}{130,30,30}
\definecolor{FdMGreen}{RGB}{25,100,60}

\setbeamercolor{structure}{fg=FdMBlue}
\setbeamercolor{title}{fg=white,bg=FdMBlue}
\setbeamercolor{frametitle}{fg=white,bg=FdMBlue}
\setbeamercolor{block title}{fg=white,bg=FdMBlue}
\setbeamercolor{block body}{bg=blue!3}
\setbeamercolor{alerted text}{fg=FdMRed}

\setmainfont{Latin Modern Roman}
\setsansfont{Latin Modern Sans}
\setmonofont{Latin Modern Mono}

\newcommand{\N}{\mathbb{N}}
\newcommand{\Z}{\mathbb{Z}}
\newcommand{\Q}{\mathbb{Q}}
\newcommand{\R}{\mathbb{R}}
\newcommand{\C}{\mathbb{C}}
\newcommand{\Pow}{\mathcal{P}}
\newcommand{\id}{\operatorname{id}}
\newcommand{\mcd}{\operatorname{mcd}}
\newcommand{\mcm}{\operatorname{mcm}}

\title[FdM I -- Clase 06]{Conjuntos I}
\subtitle{Pertenencia, inclusión y operaciones básicas}
\author{Antonio Falcó Montesinos}
\institute{Universidad CEU Cardenal Herrera}
\date{Fundamentos de las Matemáticas I}

\begin{document}

\begin{frame}
  \titlepage
\end{frame}

\begin{frame}{Objetivos de la clase}
\begin{itemize}
  \item Definir conjuntos, elementos, pertenencia e inclusión.
  \item Trabajar con unión, intersección, diferencia y complementario.
  \item Relacionar operaciones de conjuntos con conectores lógicos.
\end{itemize}
\end{frame}

\begin{frame}{Mapa de la clase}
\tableofcontents
\end{frame}

\section{Conjuntos y elementos}

\begin{frame}{Conjuntos y elementos}
\begin{block}{Pertenencia}
\begin{itemize}
  \item Un conjunto agrupa objetos llamados elementos.
  \item \(x\in A\) significa que \(x\) pertenece al conjunto \(A\).
  \item \(x\notin A\) significa que \(x\) no pertenece a \(A\).
\end{itemize}
\end{block}
\end{frame}

\begin{frame}{Conjuntos y elementos}
\begin{block}{Inclusión}
\begin{itemize}
  \item \(A\subseteq B\) significa que todo elemento de \(A\) pertenece a \(B\).
  \item Para probar una inclusión, se toma \(x\in A\) y se demuestra \(x\in B\).
  \item Para refutarla, se busca \(x\in A\) con \(x\notin B\).
\end{itemize}
\end{block}
\end{frame}

\begin{frame}{Conjuntos y elementos}
\begin{block}{Conjunto vacío}
\begin{itemize}
  \item \(\varnothing\) es el conjunto que no tiene elementos.
  \item Se cumple \(\varnothing\subseteq A\) para todo conjunto \(A\).
  \item No debe confundirse \(\varnothing\) con \(\{\varnothing\}\).
\end{itemize}
\end{block}
\end{frame}

\section{Operaciones}

\begin{frame}{Operaciones}
\begin{block}{Unión e intersección}
\begin{itemize}
  \item \(A\cup B=\{x:x\in A\text{ o }x\in B\}\).
  \item \(A\cap B=\{x:x\in A\text{ y }x\in B\}\).
  \item La unión corresponde a ``o''; la intersección corresponde a ``y''.
\end{itemize}
\end{block}
\end{frame}

\begin{frame}{Operaciones}
\begin{block}{Diferencia y complementario}
\begin{itemize}
  \item \(A\setminus B=\{x:x\in A\text{ y }x\notin B\}\).
  \item El complementario necesita un conjunto universal fijado.
  \item Sin universo, el complementario no está completamente definido.
\end{itemize}
\end{block}
\end{frame}

\section{Profundización}

\begin{frame}{Profundización}
\begin{block}{Demostrar una inclusión}
\begin{itemize}
  \item Para probar \(A\subseteq B\), se toma un elemento arbitrario \(x\in A\).
  \item Se usa la definición de \(A\) para obtener información sobre \(x\).
  \item Se prueba que \(x\in B\).
\end{itemize}
\end{block}
\end{frame}

\begin{frame}{Profundización}
\begin{block}{Pertenencia frente a inclusión}
\begin{itemize}
  \item \(x\in A\) compara un objeto con un conjunto.
  \item \(A\subseteq B\) compara dos conjuntos.
  \item Confundir \(\in\) y \(\subseteq\) suele producir enunciados sin sentido.
\end{itemize}
\end{block}
\end{frame}

\begin{frame}{Profundización}
\begin{block}{Conectores y conjuntos}
\begin{itemize}
  \item \(x\in A\cap B\) equivale a \(x\in A\) y \(x\in B\).
  \item \(x\in A\cup B\) equivale a \(x\in A\) o \(x\in B\).
  \item \(x\in A\setminus B\) equivale a \(x\in A\) y \(x\notin B\).
\end{itemize}
\end{block}
\end{frame}

\begin{frame}{Cierre}
\begin{itemize}
  \item La pertenencia habla de elementos; la inclusión habla de conjuntos.
  \item Las operaciones conjuntistas traducen conectores lógicos.
  \item La próxima clase se centra en producto cartesiano, partes y familias.
\end{itemize}
\end{frame}

\begin{frame}{Trabajo recomendado}
\begin{itemize}
  \item Releer las secciones correspondientes del manual.
  \item Identificar definiciones, símbolos nuevos y enunciados principales.
  \item Preparar dudas y ejemplos para el seminario asociado.
\end{itemize}
\end{frame}

\end{document}
